\par\Needspace{7\baselineskip}\bigskip
\begin{center}
\pdfbookmark[2]{Gauss to Olbers, Göttingen, 8 April 1813}{pallas-brief-18130408}
\textsc{Gauss} to \textsc{Olbers}, Göttingen, 8 April 1813.
\end{center}
\SourceLeaf{428}\noindent\hspace*{\parindent}. . . . .\quad I have made quite some progress with the perturbations of Pallas. Only one element remains, the epoch. For the perihelion there are (from {\AstroSymbols ♃}) 183 equations that exceed \(0''{,}1\). Once the epoch is dealt with as well, I may perhaps still undertake at least a \textls[180]{trial} set of tables. 800 equations are admittedly daunting, but the circumstance that \(7\text{\AstroSymbols ⚴}=18\text{\AstroSymbols ♃}\) (as closely as possible) allows a considerable shortcut. In the meantime Mr \textsc{Nicolai} will probably finish the perturbations by {\AstroSymbols ♄}, and perhaps another very able pupil of mine, Mr \textsc{Encke}, will undertake the perturbations by {\AstroSymbols ♂}. If your time permits, please also help look for {\AstroSymbols ⚴} soon.\quad . . . . .
\par\Needspace{7\baselineskip}\bigskip
\begin{center}
\pdfbookmark[2]{Gauss to Olbers, Göttingen, 2 July 1813}{pallas-brief-18130702}
\textsc{Gauss} to \textsc{Olbers}, Göttingen, 2 July 1813.
\end{center}
\noindent\hspace*{\parindent}. . . . .\quad The perturbations of Pallas by Jupiter, so far as I intended to calculate them, are now mostly finished. Much of the numerical calculation for the oppositions to date, however, is still outstanding.\quad . . . . .
\par
